\section{Magic of Non-Simply Connected Subsystems}
The reduced density matrix of a subsystem with a hole inside takes the following form
\[
    \rho_{A}
    =
    \frac{1}{2^{N^A - N^A_p - 2g}}
    \hat{\Pi}_A
     \hat{\Pi}_{\mathcal{L}}
    \sum_{[\mathfrak{D}]\subset A}
    \lr{
    \frac{1}{2^g}
    \sum_{\lrg{\mathfrak{H}}}
        \mathfrak{C}_{\mathfrak{D}\ominus\lrg{\mathfrak{H}}} \Sigma_{\mathfrak{D}\ominus\lrg{\mathfrak{H}}} 
    }.
\] 
The summation of the expectation value of Pauli strings to the power four is then
\[
    \frac{1}{2^{N^A}}
    \sum_{\mathsf{S}\in\tilde{\mathcal{P}}_{N^A}}
    \lrm{\Tr_A\lr{\rho_A\mathsf{S}}}^{4}
    =
    \frac{1}{2^{N^A - N_p^A - 2g}}
    \sum_{[\mathfrak{D}] \subset A}
    \lrm{
        \frac{1}{2^g}
        \sum_{\lrg{\mathfrak{H}}} \lr{\frac{\mathfrak{C}_{\mathfrak{\mathfrak{D}\ominus\lrg{\mathfrak{H}}}}}{\mathfrak{C}_\mathfrak{D}}}^4
    }
    \mathfrak{C}_\mathfrak{D}^4
\]
Because $\mathfrak{C}_\mathfrak{D}^2 \geq \mathfrak{C}_{\mathfrak{D}\ominus\lrg{\mathfrak{H}}}^2$, the series in the square bracket is less than unity,
\[
\frac{1}{2^g}
\sum_{\lrg{\mathfrak{H}}} \lr{\frac{\mathfrak{C}_{\mathfrak{\mathfrak{D}\ominus\lrg{\mathfrak{H}}}}}{\mathfrak{C}_\mathfrak{D}}}^4
\leq
\frac{1}{2^g}
\sum_{n=1}^{2^g}
=
1.
\]
Therefore, the magic part of the subsystem 2-SRE is greater than that computed by the Majorana state with the topological boundary correction,
\[
    -\ln\lrm{
        \frac{1}{2^{N^A}}
        \sum_{\mathsf{S}\in\tilde{\mathcal{P}}_{N^A}}
        \lrm{\Tr_A\lr{\rho_A\mathsf{S}}}^{4}
    }
    \geq
    -\ln\lr{
        \frac{1}{2^{N^A - N_p^A - 2g}}
        \sum_{[\mathfrak{D}] \subset A}
        \mathfrak{C}_\mathfrak{D}^4
    }.
\]
% As a result, the subsystem SRE computed using the physical spin RDM $\rho_A$ is greater than that computed using the equivalent Majorana fermionic RDM $\rho_A^F$.
% \begin{equation}\label{eq:RDM-SRE-g}
%     \tilde{M}_2\lr{\rho_A} 
%     \geq 
%     -\ln\lr{
%         \frac{
%         \sum_{[\mathfrak{D}]\subset A} \mathfrak{C}_\mathfrak{D}^4 
%         }
%         {
%         \sum_{[\mathfrak{D}]\subset A} \eta_\mathfrak{D} \mathfrak{C}_\mathfrak{D}^2
%         } 
%     }
%     =
%     \tilde{M}_2\lr{\rho_A^F}
%     .
% \end{equation}

